%% To use: Open the IJPRAI template on Overleaf at
%% https://www.overleaf.com/latex/templates/international-journal-of-pattern-recognition-and-artificial-intelligence-ijprai/nswmytycsnvn
%% Replace the template content with everything below.
%% The ws-ijprai.cls file is already in the Overleaf project.

\documentclass{ws-ijprai}

\begin{document}

\markboth{A. M. Brilliant}
{Steelmanning LLM Review: Severity Calibration Through Adversarial Self-Challenge}

\title{Steelmanning LLM Review: Severity Calibration
Through Adversarial Self-Challenge}

\author{Andrew Michael Brilliant}
\address{Applied Dynamics Research, Sapporo, Japan\\
ORCID: 0009-0004-8024-5442\\
andrew.brilliant@gmail.com}

\maketitle

\begin{abstract}
Large language models can identify methodological concerns in scientific
manuscripts but cannot reliably distinguish fatal errors from routine
limitations. We formalize this as a severity calibration problem and
present a multi-model architecture (graduated dissent) that addresses it
through adversarial steelman exchange: separated model instances argue
against their own severity ratings before an arbiter renders judgment.

On the SPOT scientific error detection benchmark ($n=50$ text-detectable
papers, externally validated ground truth), the steelman exchange achieves
36\% pass@1 with 11.7\% retraction-worthy precision, nearly double the
precision of alternative architectures, while producing the fewest top-tier
flags per paper. On a paired benchmark of 10 retracted papers with 19
non-retracted controls (author-scored with independent multi-model audit),
the steelman exchange produces fewer false retraction-worthy classifications
on controls (2/19 under independent out-of-family audit) than multi-model
pooling (3/19). On SPOT, multi-model pooling improves detection modestly
but does not improve severity calibration over rubric-based single-model
review. Both observations are consistent with an information-theoretic
bound showing that correlated evaluators cannot accumulate independent
evidence regardless of ensemble size.

All code, data, model outputs, and a replication harness are released.
\end{abstract}

\keywords{LLM evaluation; severity calibration; steelman exchange;
adversarial self-challenge; multi-model review; error detection.}


\section{Introduction}
\label{sec:intro}

A system that produces 15 findings per paper, of which 7\% are genuine,
does not perform triage; it generates a reading list with a 93\% noise floor.
The human reviewer must manually evaluate every finding, because the system
provides no signal about which findings matter. This is the current
state of LLM-based scientific manuscript review.

The practical bottleneck is no longer detection alone. Current frontier models
can identify methodological concerns in scientific manuscripts at rates
comparable to individual human reviewers.\cite{liang2024feedback} The SPOT
benchmark\cite{son2025spot} measures whether a model finds a documented error
somewhere in its output; the FLAWS benchmark\cite{xi2025flaws}
measures error localization. Detection performance continues to improve, but
unranked findings are now the practical bottleneck.

The problem is \emph{severity calibration}: the inability to distinguish fatal
errors from routine limitations. A paper with a catastrophic statistical error
receives the same diplomatic review format as a paper with minor reporting
gaps (several strengths, several concerns, a hedged
conclusion).\cite{brilliant2026selfcorrection} This creates a practical
bottleneck: as AI-assisted writing increases submission volumes while making
surface quality less diagnostic,\cite{kobak2025llm} editorial workflows need
automated triage that can rank problems by severity, not merely list them.

We present three contributions:

\begin{enumerate}
\item A multi-model architecture in which separated model instances
exchange steelman arguments challenging each other's severity ratings,
converting undifferentiated findings into graded severity labels
(Section~\ref{sec:method}).

\item Empirical evidence from two independent benchmarks showing that
the steelman exchange improves detection and nearly doubles within-tier
retraction-worthy precision compared to alternative architectures, while
multi-model pooling without adversarial exchange provides no severity
calibration improvement over rubric-based single-model review
(Sections~\ref{sec:spot_results}--\ref{sec:retracted_results}).

\item An information-theoretic motivation explaining why pooling
correlated evaluators should not be expected to accumulate independent severity
evidence, and why adversarial exchange is a structurally plausible remedy
(Section~\ref{sec:theory}).
\end{enumerate}

Table~\ref{tab:headline} summarizes the empirical results.

\begin{table}[h]
\tbl{Results across two benchmarks. B3 and graduated dissent (GD) use the
same models and arbiter; they differ only in the steelman exchange.
Retracted-paper results use out-of-family audit (Grok~4, Mistral Large~2,
disjoint from evaluation pipeline).}
{\begin{tabular}{@{}lccccc@{}}
\toprule
& \multicolumn{2}{c}{Retracted Papers ($n=10$)} &
\multicolumn{3}{c}{SPOT ($n=50$)} \\
\cline{2-3} \cline{4-6}
Condition & Detection$^\dagger$ & FP Rate & pass@1 & RW-prec. & RW-yield \\
\colrule
B1: Single model & 6/10 & 0/19 & 24\% & n/a & n/a \\
B2: Single + rubric & 9/10 & 0/19 & 28\% & 7.0\% & 20\% \\
B3: Multi-model pool & 7/10 & 3/19 & 30\% & 6.2\% & 20\% \\
GD: Steelman & 10/10 & 2/19 & 36\% & 11.7\% & 32\% \\
\botrule
\end{tabular}\label{tab:headline}}
\begin{tabnote}
$^\dagger$Out-of-family audit scoring. In-family audit produced 0/19 FP
for GD and 2/19 for B3; independent verification revised GD to 2/19 and B3
to 3/19.
\end{tabnote}
\end{table}


\section{Method: Multi-Model Evaluation with Steelman Exchange}
\label{sec:method}

The graduated dissent protocol\cite{brilliant2026graduated} treats evaluation
as a multi-stage adversarial process.

\subsection{Architecture}

Two \textbf{provers} independently review an anonymized manuscript in separated
contexts, rating each finding's severity (retraction-worthy, major-revision,
minor). A \textbf{steelman exchange} then requires each prover to construct the
strongest case against their own severity ratings, arguing that their own
findings are less severe than initially assessed. An \textbf{arbiter} receives
all material and applies a conservative severity rubric.

\subsection{Severity rubric}

Every prompt includes the following definitions:

\begin{quote}
\textbf{Retraction-worthy:} The error means the paper's central conclusions
cannot be supported by the data as presented. Not ``could be improved'' but
``fundamentally broken.''

\textbf{Major-revision:} A real concern that could change conclusions if
addressed differently, but does not definitively invalidate them.

\textbf{Minor:} Valid criticism applicable to most published papers. Would not
change conclusions.
\end{quote}

\subsection{Why the steelman exchange targets severity}

Standard multi-model pooling (B3 in our experiments) combines findings from two
models but does not challenge severity ratings. If both models share the
tendency to over-rate severity on certain error types (a well-documented
consequence of RLHF training\cite{brilliant2026selfcorrection}), pooling
reinforces the overestimate.

The steelman exchange disrupts this by requiring each model to argue that its
own severity rating is too high. This activates a different prediction target
than the default review framing, partially decorrelating the severity assessment
from the shared failure modes that produce correlated over-rating.

\subsection{Models}

\begin{table}[h]
\tbl{Models used across both benchmarks. Cross-family diversity provides
partial decorrelation of failure modes.\cite{chen2024reconcile}}
{\begin{tabular}{@{}llll@{}}
\toprule
Role & Model & Provider & Family \\
\colrule
Prover A & GPT-5.4 & OpenAI & GPT \\
Prover B & DeepSeek V3.2 & DeepSeek & DeepSeek \\
Arbiter & Claude Opus 4.6 & Anthropic & Claude \\
\botrule
\end{tabular}\label{tab:models}}
\end{table}

\subsection{Budgeted escalation and termination}

The protocol includes a natural termination condition: when inter-prover
agreement exceeds a threshold ($\theta_{\text{accept}} = 0.90$), the steelman
exchange is skipped and the arbiter receives only the convergent reviews.

In the SPOT evaluation ($n=50$), the escalation gate fired correctly: 3 papers
hit L0 (agreement $\geq 0.90$) and skipped the steelman exchange; the
remaining 47 escalated to L2 and ran the full adversarial exchange. This
provides a stop condition that RLHF-trained review lacks: the protocol
terminates when adversarial challenge produces no surviving severity disputes.


\section{Benchmark 1: SPOT (Detection and Severity Ranking)}
\label{sec:spot_results}

\subsection{Design}

We evaluated all four conditions on the text-detectable subset of the SPOT
benchmark,\cite{son2025spot} which contains real published papers with
human-validated errors. Papers with figure-dependent errors were excluded (our
pipeline is text-only), yielding $n=50$ papers. All papers were anonymized
before evaluation.

\subsection{Detection results}

GD achieves 36\% pass@1, the highest among evaluated conditions. Multi-model
pooling (B3) provides modest improvement over single-model review (+6pp over
B1), and the steelman exchange provides a further +6pp beyond pooling.
Because our evaluation uses a text-detectable subset ($n=50$) with different
prompting, while the published o3 result used the full multimodal benchmark
($n=83$), the comparison with o3 should be interpreted as descriptive context
rather than a controlled comparison.

\subsection{Severity ranking: the actionability argument}

Pass@1 measures whether the system found the annotated error somewhere in
its output, typically a list of 10--15 findings per paper. For triage, the
relevant question is not ``did it find something?'' but ``did it tell you
this one is the problem?''

GD finds more errors (19 TPs vs.\ B3's 16) and grades them better: the most
true positives in the retraction-worthy tier (9 vs.\ B3's 6), the fewest
total RW flags (77 vs.\ 97), and nearly double the within-tier precision
(11.7\% vs.\ 6.2\%). A reviewer who reads only retraction-worthy findings
gets the most genuine errors in the fewest flags.

\subsection{Budget-matched controls}

To control for compute budget, two additional conditions were evaluated at
the same five API calls as GD ($n=49$; one paper excluded due to API content
filter).

\begin{table}[h]
\tbl{Budget-matched conditions on SPOT ($n=49$). All three use five
API calls with different framing for the additional passes.}
{\begin{tabular}{@{}lccccc@{}}
\toprule
Condition & Calls & pass@1 & RW-prec. & RW-yield & RW/paper \\
\colrule
B3+: Neutral reflect & 5 & 22.4\% & 7.2\% (6/83) & 19.4\% & 1.69 \\
B3++: Defend ratings & 5 & 36.7\% & 6.2\% (7/112) & 19.4\% & 2.29 \\
GD: Steelman & 5 & 34.7\% & 10.5\% (8/76) & 29.6\% & 1.55 \\
B3 (3-call ref.) & 3 & 28.6\% & 5.2\% (5/96) & 17.2\% & 1.96 \\
\botrule
\end{tabular}\label{tab:budget}}
\end{table}

\subsection{Scoring validation}

True positive verification uses a five-layer validation stack:
(1)~SPOT's author-acknowledged ground truth;
(2)~SPOT's own LLM-as-judge matching protocol applied verbatim;
(3)~two adversarial model audits following an 8-step verification protocol;
(4)~manual human spot-check of a random sample (5/19 GD TPs, all confirmed);
and (5)~blinded human verification by a graduate student who independently
re-scored all final true-positive decisions without access to condition
labels, achieving 100\% agreement.


\section{Benchmark 2: Retracted Papers (Severity Calibration)}
\label{sec:retracted_results}

\subsection{Design}

The SPOT benchmark tests detection and severity ranking but does not test
false positives on sound papers. We constructed a benchmark of 10 retracted
papers with documented methodological errors, paired with 19 non-retracted
controls. All retracted papers post-date model training cutoffs. Ground truth
was encoded as keyword-decomposed retraction causes. Full details are
in Ref.~\refcite{brilliant2026benchmark}.

Detection and false-positive scoring was performed by the first author,
verified by in-family audit (1,395 judgments) and out-of-family audit (1,318
findings scored by Grok~4 and Mistral Large~2, disjoint from the pipeline).

\subsection{Results}

Under OOF audit, GD detects all 10 retraction-causing errors (10/10) versus
B3's 7/10, and produces fewer false positives (2/19 vs.\ 3/19). The original
0/19 FP result under in-family audit does not survive independent verification.

The steelman exchange's contribution on this benchmark is not detection but
severity calibration: producing fewer false retraction-worthy classifications
on controls than multi-model pooling.


\section{Cross-Benchmark Analysis}

The two benchmarks test different capabilities and tell a consistent story.
On SPOT, the steelman exchange's primary effect is grading quality. On
retracted papers under independent audit, the primary effect is
discrimination. Both serve the same practical function: converting raw model
output into ranked severity labels.

The limited effect of multi-model pooling (B3) is equally informative. On
SPOT, B3 provides modest detection improvement but no severity calibration
improvement: its RW-precision (6.2\%) and RW-yield (20\%) match rubric-based
single-model B2.


\section{Why Pooling Does Not Improve Severity Calibration}
\label{sec:theory}

The empirical pattern admits a formal explanation. This section provides
information-theoretic bounds that motivate the observations.

\subsection{Setup}

Let $G \sim p(g \mid x)$ be a generator, $S \sim p(s \mid x, g)$ a selector,
and $T := \mathbb{I}\{G = H^\star(X)\}$ the correctness indicator.

\textbf{Modeling assumption.} We assume $S \perp T \mid (G, Z)$ where
$Z \in \mathcal{Z}$ indexes shared failure modes.

\begin{theorem}[Information Bound via Shared Blind Spots]
\label{thm:info_bound}
Assume $S \perp T \mid (G, Z)$. Then:
$I(T; S \mid G) \leq I(T; Z \mid G)$.
If $T \perp Z \mid G$, then $I(T; S \mid G) = 0$.
\end{theorem}

\begin{lemma}[Repeated Self-Critique Bound]
\label{lem:repeated}
For $k$ selectors with $(S_1, \ldots, S_k) \perp T \mid (G, Z)$ jointly:
$I(T; A_{1:k} \mid G) \leq I(T; Z \mid G)$.
\end{lemma}

Under correlated error, $k$ evaluations accumulate subjective confidence
without accumulating objective evidence. This is consistent with B3's failure
to improve severity calibration on either benchmark.

\begin{proposition}[Sufficient Conditions for Positive Evidence]
\label{prop:positive}
Assume a selector satisfies $\Pr(A=1 \mid T=0) \leq 1-\epsilon$ and
$\Pr(A=1 \mid T=1) \geq \delta$ for some $\delta > 1-\epsilon > 0$. Then
acceptance provides $\log_2 \frac{\delta}{1-\epsilon} > 0$ bits of evidence
about correctness.
\end{proposition}

The steelman exchange is designed to partially satisfy this condition,
reducing false acceptance of overblown severity classifications.


\section{Related Work}

Liang et al.\cite{liang2024feedback} found GPT-4 feedback overlapped with
human reviewers at 30--39\%. SPOT\cite{son2025spot} tests detection of
verified errors (o3 achieves 18.4\% pass@1).
FLAWS\cite{xi2025flaws} tests error localization.
Pub-Guard-LLM\cite{chen2025pubguard} achieves $\sim$91\% on binary retraction
classification. None test severity calibration.

Du et al.\cite{du2024debate} showed multi-agent debate improves factuality.
Huang et al.\cite{huang2024selfcorrect} showed debate performs no better than
self-consistency at matched compute. Chen et al.\cite{chen2024reconcile} found
model family diversity contributes 6.8\% of gains. Our architecture differs:
models argue against their own positions (steelman), target severity rather
than correctness, and preserve disagreement.

Kim et al.\cite{kim2025correlated} document error correlation across 350+
models. Li et al.\cite{li2024biases} identified 12 latent biases.
Tsui\cite{tsui2025blindspot} found a 64.5\% self-correction failure rate.
These motivate the adversarial self-challenge approach.


\section{Limitations}

\begin{enumerate}
\item \textbf{Sample size.} The retracted-paper benchmark comprises $n=10$
retracted papers with 19 controls. These results motivate larger-scale
replication.

\item \textbf{SPOT subset.} Results are on the text-detectable subset
($n=50$), excluding $\sim$30 papers with figure-dependent errors.

\item \textbf{Disputed scoring.} One SPOT discordant pair (paper\_014) was
classified differently by two independent auditors. Strict scoring is reported
throughout.

\item \textbf{Retracted-paper scoring.} Initial scoring was performed by the
first author, with independent multi-model audit and out-of-family audit
verification. The in-family 0/19 FP result does not survive OOF verification
(2/19). All audit data are released for independent adjudication.

\item \textbf{Theoretical bounds require falsifiable Z.} The information
bound is non-vacuous only when $Z$ is independently constrained.

\item \textbf{Context separation is heuristic.} The steelman exchange
partially disrupts error correlation but does not break parameter-level
correlations from shared training.

\item \textbf{Cost.} The full pipeline costs approximately \$1.25 per paper.
\end{enumerate}


\section{Discussion}

All four architectures exceed SPOT's published o3 baseline on text-detectable
errors. Detection performance has improved enough that unranked findings are
now the practical bottleneck.

The steelman exchange works not by finding more errors but by generating
better judgments about errors already found. This is consistent with the
theoretical framework: pooling correlated evaluators reinforces shared biases
(Lemma~\ref{lem:repeated}), while adversarial challenge partially disrupts
the correlation (Proposition~\ref{prop:positive}).

To support practical adoption, the complete pipeline (including prompts,
scoring code, model configurations, and a browser-based interface) is
released as open-source infrastructure.


\section{Conclusion}

Actionable severity ranking, not raw detection, is the bottleneck for
practical LLM-based scientific evaluation. Among the evaluated conditions,
the steelman exchange is the only intervention that improves severity
calibration across both benchmarks.

On SPOT ($n=50$), it achieves 36\% pass@1 with 11.7\% retraction-worthy
precision and the fewest top-tier flags per paper. On retracted papers
($n=10$), GD detects all 10 retraction-causing errors with 2/19 false
positives on controls under independent out-of-family audit, compared to
B3's 7/10 detection with 3/19 false positives.

All code, data, model outputs, and a replication harness are available at
\texttt{https://github.com/AndBrilliant/GraduatedDissentBench}.


\section*{Data Availability}
All code, model outputs, scoring data, and audit reports are publicly
available at \texttt{https://github.com/AndBrilliant/GraduatedDissentBench}.

\section*{Use of AI Tools}
This research evaluates large language models as the subject of study.
All models used, their roles, and their outputs are documented in the
methodology and released in the repository.

\section*{Conflicts of Interest}
The author is the developer of the graduated dissent protocol and the
associated open-source software. The author has no financial interest in the
models or API services used. No external funding was received.


\bibliographystyle{ws-ijprai}
\begin{thebibliography}{99}

\bibitem{brilliant2026selfcorrection}
A.~M.~Brilliant, Limits of self-correction in LLMs: An information-theoretic
analysis of correlated errors, \textit{Preprints.org}, DOI:
10.20944/preprints202601.0892.v3 (2026).

\bibitem{brilliant2026graduated}
A.~M.~Brilliant, Graduated dissent: Budgeted disagreement resolution for
multi-model inference, \textit{Preprints.org}, DOI:
10.20944/preprints202603.1830.v1 (2026).

\bibitem{brilliant2026benchmark}
A.~M.~Brilliant, The graduated dissent benchmark: A paired-ablation resource
for multi-model manuscript review, in preparation (2026).

\bibitem{huang2024selfcorrect}
J.~Huang et al., Large language models cannot self-correct reasoning yet,
\textit{Proc. ICLR} (2024).

\bibitem{tsui2025blindspot}
K.~Tsui, Self-correction bench, \textit{arXiv:2507.02778} (2025).

\bibitem{kim2025correlated}
E.~Kim, A.~Garg, K.~Peng and N.~Garg, Correlated errors in large language
models, \textit{Proc. ICML} (2025).

\bibitem{son2025spot}
G.~Son et al., SPOT: A benchmark for automated verification of scientific
research, \textit{arXiv:2505.11855} (2025).

\bibitem{xi2025flaws}
Y.~Xi et al., FLAWS: A benchmark for error identification,
\textit{arXiv:2511.21843} (2025).

\bibitem{chen2025pubguard}
X.~Chen et al., Pub-Guard-LLM: Detecting retracted biomedical articles,
\textit{arXiv:2502.15429} (2025).

\bibitem{du2024debate}
Y.~Du et al., Improving factuality through multiagent debate,
\textit{Proc. ICML} (2024).

\bibitem{chen2024reconcile}
J.~C.-Y.~Chen, S.~Saha and M.~Bansal, ReConcile: Round-table conference
improves reasoning, \textit{Proc. ACL} (2024).

\bibitem{liang2024feedback}
W.~Liang et al., Can large language models provide useful feedback on
research papers? \textit{NEJM AI} \textbf{1}(8) (2024).

\bibitem{li2024biases}
J.~Li et al., Justice or prejudice? Quantifying biases in LLM-as-a-Judge,
\textit{arXiv:2410.02736} (2024).

\bibitem{wataoka2024selfpref}
K.~Wataoka, T.~Takahashi and R.~Ri, Self-preference bias in LLM-as-a-Judge,
\textit{arXiv:2410.21819} (2024).

\bibitem{lightman2023verify}
H.~Lightman et al., Let's verify step by step, \textit{arXiv:2305.20050}
(2023).

\bibitem{baba2025prover}
K.~Baba et al., Prover Agent for formal mathematical proofs,
\textit{arXiv:2506.19923} (2025).

\bibitem{kobak2025llm}
D.~Kobak et al., LLM-assisted writing in biomedical publications,
\textit{Science Advances} \textbf{11}, eadt3813 (2025).

\end{thebibliography}

\end{document}
